\section{Limitations}
\label{sec:limitations}

Synthetic agents are not real depositors, and the behavioral calibration to human decision-making is necessarily partial. The LLM agents in this study are prompted to simulate depositor reasoning, but they do not replicate the full emotional, cognitive, and social context of a human facing an actual bank-run scenario. In particular, loss aversion, social network effects, and media consumption patterns are compressed into the information-environment parameterization rather than modeled at individual agent level. The quantitative $\Pr(\text{run})$ estimates should therefore be interpreted as characterizations of the synthetic lab environment rather than as predictions of real-world run probabilities in any specific institution or market.

The analysis maintains a single-bank focus throughout; the Acemoglu network extension in Corollary~\ref{cor:network} is a theoretical result and is not directly tested in the E01--E08 experimental battery. Multi-bank cascade dynamics, where a synchronized withdrawal at one institution propagates through the interbank exposure graph $G$ to trigger secondary runs at connected institutions, are left for future work. The empirical phase diagram in Figure~\ref{fig:e03} should not be applied to multi-bank systemic risk assessment without a validated extension of the model to the network setting.

Foundation-model evolution represents a significant threat to longitudinal replicability. The LLM architectures available at the time of writing (GPT-4, Claude 3, Llama-3, Mistral-7B) will be superseded by models with different behavioral profiles, and there is no guarantee that the amplification factor $\lambda$ estimated from E03 will generalize to future model generations. We mitigate this concern by pinning all simulation runs to specific model version identifiers recorded in the pre-registration and by providing complete prompt logs in the supplementary material, enabling partial replication with future model versions. However, any regulatory or policy recommendation derived from this study should be treated as requiring periodic re-calibration as the foundation-model landscape evolves.

Prompt sensitivity is mitigated by the E06 robustness audit (Section~\ref{sec:robust}) but not eliminated. The five prompt variants tested cover a range of syntactic framings but do not exhaust the space of possible prompts; a sufficiently adversarial prompt design could plausibly reverse the main effect. We report the range of effect sizes across E06 variants as a transparency measure, and we note that the pre-registration commits to treating any reversal in E06 as disconfirming evidence. Nevertheless, the findings are conditional on the set of prompts tested, and replication with independently constructed prompts by researchers not affiliated with this study would provide stronger evidence of generalizability.

% TODO: add a limitation paragraph on the absence of regulatory feedback loops: the current model treats regulatory intervention as exogenous (E08), but in reality regulators observe and respond to agent behavior, creating a co-evolutionary dynamic not captured here.
